\chapter{Feature Engineering}
\label{ch:feature-engineering}

In deze fase van het onderzoek werd feature engineering toegepast om op basis van ruwe wedstrijddata nieuwe, informatieve variabelen te construeren. 
Het doel van deze stap is het verhogen van de voorspellende kracht van de dataset en het toevoegen 
van contextuele en interpreteerbare informatie voor de machine learning-modellen. Alle dynamische features werden zodanig geconstrueerd 
dat uitsluitend gebruik werd gemaakt van historische informatie om datalekage te vermijden.

\vspace{5pt}

\subsection{Seizoensvariabele}

Elke wedstrijd werd toegewezen aan een specifiek seizoen op basis van de datum. 
Wedstrijden gespeeld tussen juli en december werden toegewezen aan het seizoen startjaar/startjaar+1, 
terwijl wedstrijden tussen januari en juni werden toegewezen aan het vorige startjaar/startjaar+1 seizoen. 
Deze variabele maakt het mogelijk om prestaties te groeperen binnen competitieve contexten.

\vspace{5pt}

\subsection{Dynamische rangschikking}

Voor elke wedstrijd werd een tussentijdse rangschikking berekend op basis van alle voorgaande wedstrijden binnen hetzelfde seizoen. 
Teams kregen 3 punten voor een overwinning, 1 punt voor een gelijkspel en 0 punten voor een nederlaag. 
Daarnaast werd het doelsaldo berekend als het verschil tussen gemaakte en geïncasseerde doelpunten.

De rangschikking werd dynamisch bepaald op elk wedstrijdmoment door teams te sorteren op:
\begin{itemize}
    \item totaal aantal punten (aflopend)
    \item doelsaldo (aflopend)
\end{itemize}

Deze aanpak simuleert de evolutie van de competitietabel doorheen het seizoen.

\vspace{5pt}

\subsection{Recente vorm}

De vorm van een team werd gedefinieerd als het gemiddelde aantal behaalde punten over de laatste vijf gespeelde wedstrijden vóór de betreffende wedstrijd. 
Elke wedstrijdresultaat werd omgezet naar een numerieke score (3 voor winst, 1 voor gelijkspel, 0 voor verlies). 

Formeel:

\[
\text{Form}_{t} = \frac{1}{5} \sum_{i=t-5}^{t-1} P_i
\]

waarbij $P_i$ het aantal behaalde punten in wedstrijd $i$ is.

Deze variabele geeft een indicatie van de recente prestatiecurve van een team en vangt tijdelijke fluctuaties in vorm op.

\vspace{5pt}

\subsection{Head-to-head statistieken}

Voor elke wedstrijd werden historische onderlinge confrontaties tussen beide teams geanalyseerd. 
Hierbij werd gekeken naar de laatste drie onderlinge wedstrijden vóór het huidige moment. De volgende variabelen werden berekend:

\begin{itemize}
    \item aantal overwinningen van het thuisteam
    \item aantal overwinningen van het uitteam
\end{itemize}

Op basis hiervan werd een binaire classificatie afgeleid die aangeeft welk team historisch voordeel heeft in de onderlinge duels.

\vspace{5pt}

\subsection{Titel- en degradatiewedstrijden}

Wedstrijden werden geclassificeerd als titelduels indien aan volgende voorwaarden werd voldaan:
\begin{itemize}
    \item beide teams staan in de top 2 van de rangschikking
    \item beide teams hebben meer dan 40 punten
    \item het puntenverschil bedraagt maximaal 5 punten
\end{itemize}

Degradatiewedstrijden werden gedefinieerd als wedstrijden die laat in het seizoen plaatsvinden waarbij 
beide teams zich in de onderste regionen van de rangschikking bevinden (laatste drie plaatsen) en het puntenverschil maximaal 3 punten bedraagt.

\vspace{5pt}

\subsection{Rivaliteitsindicator}

Op basis van vooraf gedefinieerde clubparen werd een binaire variabele toegevoegd die aangeeft of een wedstrijd een klassieke rivaliteit of derby betreft. 
Dit omvat onder andere regionale derby’s en historisch belangrijke confrontaties.

\vspace{5pt}

\subsection{Speelstijl van teams}

Voor elk team werd een speelstijlscore berekend op basis van geaggregeerde prestatie-indicatoren over de laatste vijf wedstrijden. 
Hierbij werden onder meer het aantal doelpunten, schoten en corners meegenomen.

De stijlscore werd gedefinieerd als:

\[
\text{Style Score} = (\text{doelpunten} + 0.3 \cdot \text{schoten}) - (\text{tegendoelpunten} + 0.3 \cdot \text{schoten tegen})
\]

Op basis van de verdeling van deze score werden teams ingedeeld in drie categorieën:
\begin{itemize}
    \item Aanvallend
    \item Gebalanceerd
    \item Verdedigend
\end{itemize}

De drempels werden bepaald op basis van de 33e en 66e percentielen.

\vspace{5pt}

\subsection{Contextuele variabelen}

Daarnaast werden verschillende contextuele features toegevoegd:

\begin{itemize}
    \item Weekdag en weekendindicator
    \item Avondwedstrijden (start ≥ 18:00)
    \item Weersomstandigheden (regen, harde regen, hoge temperatuur)
    \item Stadioncapaciteit (gecategoriseerd in klein, middel en groot via kwantielen)
\end{itemize}

\vspace{5pt}

\subsection{Relatieve verschillen}

Tot slot werden relatieve verschillen tussen teams berekend:

\begin{itemize}
    \item Rangverschil: verschil in positie op de dynamische ranglijst
    \item Vormverschil: verschil in recente vorm (5 wedstrijden)
\end{itemize}

Deze features laten toe om directe comparatieve sterkte tussen teams te modelleren.

\vspace{5pt}

Door deze combinatie van domeinkennis, statistische aggregatie en contextuele informatie werd een uitgebreide feature space gecreëerd. 
Deze verrijkte dataset vormt de input voor de verdere exploratieve en statische data-analyse en de machine learning-modellen.